% Lösungen Einheit 4 – Nr. 45–49
\einheitenkopf[e4]{Einheit 4 von 4 · Sachaufgaben}

\erg{45}{a) $x_1$ \quad b) $x_2$ \quad c) $x_1$ \quad d) beide \quad e) $x_2$}
\erg{46}{a) 10 und $-10$ \quad b) 11 und $-11$ \quad c) $x^2 = 36$; 6 und $-6$ \quad d) $x^2 = 81$; 9 und $-9$ \quad e) $x^2 + 20 = 84$, $x^2 = 64$; 8 und $-8$ \quad f) $x\cdot(x + 1) = 56$, $x^2 + x - 56 = 0$; $x_1 = 7$, $x_2 = -8$; die Zahlenpaare 7 und 8 oder $-8$ und $-7$}
\erg{47}{a) $x\cdot(x + 3) = 40$ \quad b) $x^2 + 3x - 40 = 0$; $x_1 = 5$, $x_2 = -8$ \quad c) $x_1 = 5$, denn eine Breite kann nicht negativ sein. \quad d) Das Beet ist 5\,m breit und 8\,m lang. \quad e) $5\,\text{m}\cdot 8\,\text{m} = 40\,\text{m}^2$ (wA) \quad f) $(x + 2)\cdot(x + 6) = 96$; $x^2 + 8x - 84 = 0$; $-4 \pm\sqrt{100}$; $x_1 = 6$, $x_2 = -14$ (keine Länge). Das Beet ist 6\,m breit und 10\,m lang; Probe: $8\,\text{m}\cdot 12\,\text{m} = 96\,\text{m}^2$.}
\erg{48}{a) Die negative Lösung ist keine Länge. Richtig: nur $x = 6$; das Rechteck ist 6\,cm breit und 8\,cm lang. \quad b) Lisa hat mit dem Umfang statt mit dem Flächeninhalt gerechnet. Richtig: $x\cdot(x + 5) = 24$, $x^2 + 5x - 24 = 0$; $x_1 = 3$, $x_2 = -8$; das Rechteck ist 3\,m breit und 8\,m lang.}
\erg{49}{a) Längen können nicht negativ sein; die negative Lösung der Gleichung fällt weg. \quad b) $5^2 = 25$ und $(-5)^2 = 25$; die Frage schließt keine der beiden Zahlen aus.}
